\section{Forecasting methods}\label{forecasting-methods}

\subsection{4.1 Quantile loss and
coverage}\label{quantile-loss-and-coverage}

Forecast accuracy is measured by the q=0.90 pinball loss

\[
L_{0.90}(y,\widehat q)=\left(0.90-\mathbf{1}\{y<\widehat q\}\right)(y-\widehat q).
\]

Lower values are better. Marginal coverage is the share of realised
outcomes satisfying \(y\leq\widehat q\). A well-calibrated q90 forecast
should have coverage near 0.90, but coverage alone is insufficient
because arbitrarily high forecasts can over-cover. Pinball loss jointly
penalises underprediction and excessive conservatism.

\subsection{4.2 Frozen base quantile
model}\label{frozen-base-quantile-model}

A quantile LightGBM model is trained through December 2020 using
seasonal terms, recent item-level price changes, backward three-month
inflation, rolling means and volatility, and cross-sectional food-price
summaries. Categorical identities describe items, markets or entities
where available. The model uses the quantile objective at
\(\alpha=0.90\), learning rate 0.05, 31 leaves, 250 minimum observations
per leaf, feature and bagging fractions of 0.85, L2 penalty 2, and 120
boosting rounds. The base model is frozen before 2021 so that all
subsequent residuals are honest pseudo-out-of-sample errors. LightGBM is
used as a flexible cross-sectional ranking model, not as the sole
adaptation mechanism (\citeproc{ref-ke2017lightgbm}{Ke et al. 2017}).

\subsection{4.3 Four expert forecasts}\label{four-expert-forecasts}

\textbf{Table 2. Expert forecasts supplied to the controller}

\begin{longtable}[]{@{}
  >{\raggedright\arraybackslash}p{(\columnwidth - 4\tabcolsep) * \real{0.3243}}
  >{\raggedright\arraybackslash}p{(\columnwidth - 4\tabcolsep) * \real{0.3243}}
  >{\raggedright\arraybackslash}p{(\columnwidth - 4\tabcolsep) * \real{0.3243}}@{}}
\toprule\noalign{}
\begin{minipage}[b]{\linewidth}\raggedright
\textbf{Expert}
\end{minipage} & \begin{minipage}[b]{\linewidth}\raggedright
\textbf{Definition}
\end{minipage} & \begin{minipage}[b]{\linewidth}\raggedright
\textbf{Adaptation channel}
\end{minipage} \\
\midrule\noalign{}
\endhead
\bottomrule\noalign{}
\endlastfoot
Static empirical q90 & A single historical 90th percentile estimated
before the test period. & None; stable-regime benchmark. \\
Fixed-persistence gate & Static q90 outside stress; frozen base plus
EWMA common residual shift for six months after a stress alarm. & Fast
common level-shift adaptation. \\
Global conformal (12m) & Frozen base plus the 90th percentile of matured
residuals from the previous 12 origin months. & Rolling global
calibration. \\
Pure hierarchical conformal & Frozen base plus a partially pooled
residual quantile using global and group-specific residual histories. &
Commodity/market heterogeneity. \\
\end{longtable}

Notes: All residual histories are truncated at t-3 at forecast origin t.
The conformal experts are empirical residual calibrators and do not
claim finite-sample conditional coverage under dependence and
distribution shift.

\subsection{4.4 Stress signal and fixed-persistence
expert}\label{stress-signal-and-fixed-persistence-expert}

The observable food stress signal is formed from robust positive
z-scores of the cross-sectional median and 90th percentile of current
monthly price changes and backward three-month inflation. Each source is
standardised against its pre-break history. The fixed-persistence expert
activates when the food stress percentile exceeds the 95th pre-break
percentile and remains active for six forecast origins.

When active, the forecast replaces the static quantile with the frozen
base plus an exponentially weighted estimate of the 90th percentile of
matured residuals. The EWMA update rate is 0.20. The six-month
persistence was selected before the extended evaluation and is
intentionally simple; it later becomes a source of recovery
overforecasting.

\subsection{4.5 Global and hierarchical residual
calibration}\label{global-and-hierarchical-residual-calibration}

Global conformal calibration adds the empirical 90th percentile of
residuals from the preceding 12 matured origin months to the frozen base
forecast. Hierarchical calibration uses a six-month window and shrinks
group-specific residual quantiles toward a global residual quantile with
prior strength \(\kappa=50\). In the NBS continuation, the local group
is the food item. In WFP, the hierarchy includes commodity, market, and
market-commodity-unit series effects.

The hierarchy is an empirical partial-pooling estimator. It can improve
weak-group calibration when local histories are sparse, but it can also
increase pooled loss if heterogeneous adjustments respond to noise. This
trade-off motivates the controller rather than selecting hierarchy as a
universally preferred model.

\subsection{4.6 Regime-aware
meta-controller}\label{regime-aware-meta-controller}

The controller combines the four expert forecasts. It first estimates,
from 2021 only, each expert's expected loss within six regime cells
defined by three levels of current stress and an indicator of high
residual heterogeneity. Cell means are shrunk toward the global 2021
expert-loss vector with prior strength three.

At origin t, the controller updates an exponentially weighted
expert-loss state only when the corresponding target has matured. Let
\(p_{j,r_t}\) denote the normalised 2021 prior loss of expert j in the
current regime \(r_t\), and let \(\bar\ell_{j,t}\) denote the normalised
delayed EWMA loss. The combined score is

\[
s_{jt}=\gamma p_{j,r_t}+(1-\gamma)\bar\ell_{j,t},
\]

with \(\gamma=0.75\). Target mixture weights are obtained from a softmax
transformation

\[
w^{\star}_{jt}=\frac{\exp(-\eta s_{jt})}{\sum_k\exp(-\eta s_{kt})},
\]

where \(\eta=10\). The selected inertia parameter equals one, so the
target weights are used directly after delayed-loss updating. When the
inherited fixed-persistence alarm is active, 75\% of the final weight is
assigned to the fixed-gate expert and the remaining 25\% is distributed
according to the loss-aware mixture. The controller forecast is the
convex combination of the four expert q90 forecasts.

\subsection{4.7 Hyperparameter selection and
evaluation}\label{hyperparameter-selection-and-evaluation}

Controller hyperparameters are selected on 2022 only from 432
configurations. The validation objective averages loss relative to the
static expert across NBS and WFP, adds 0.30 times the absolute q90
coverage error, and adds 0.02 times regime-switch and weight-turnover
friction. This prevents selection on one source or raw loss scale alone.

Evaluation reports pooled pinball loss, equal-month loss, coverage,
subgroup coverage dispersion, dynamic regret relative to the best expert
in each month, and the share of the static-to-oracle gap closed.
Month-block bootstrap intervals resample forecast-origin months so that
common within-month shocks remain intact. Simulations contain 60 series
in six latent groups, 96 months, Student-t disturbances, a three-month
feedback delay, and six no-break or structural-break scenarios.
